\documentclass{article}
\usepackage[margin=2cm]{geometry}
\usepackage{amsmath}
\usepackage{graphicx}
\usepackage{booktabs}
\usepackage[utf8]{inputenc}
\usepackage{ragged2e}
\setlength{\emergencystretch}{3em}
\begin{document}
\sloppy

\noindent\hfill\includegraphics[height=1.8em]{assets/figure_header_logo.png}\par\noindent\hrule\vspace{4pt}

e,  e, for all i  j. The free energy in the specific case of a vector multiplet in the
adjoint, a single antisymmetric hypermultiplet, and Ng fundamental hypermultiplets then is.


\begin{equation}
\begin{array}{c c}{{F(\lambda_{i})=\sum_{i\neq j}[F_{V}(\lambda_{i}-\lambda_{j})+F_{V}(\lambda_{i}+\lambda_{j})+F_{H}(\lambda_{i}-\lambda_{j})+F_{H}(\lambda_{i}+\lambda_{j})]}}&{{}}\\ {{}}&{{\displaystyle+\sum_{i}[F_{V}(2\lambda_{i})+F_{V}(-2\lambda_{i})+N_{f}F_{H}(\lambda_{i})+N_{f}F_{H}(-\lambda_{i})]}}&{{}}&{{}}&{{}}\end{array}
\end{equation}


The next step is to look for extrema of this function in the specific case of A,  0 for all.
i. Extrema in the case of non-positive A, can be obtained from these through action of the
Weyl group. To calculate the extrema, one first assumes that as N -> oo, the vevs scale as
X; = Nx, for  > 0 and x of order O(N). One then introduces a density function


\begin{equation}
\rho(x)=\frac{1}{N}\sum_{i=1}^{N}\delta(x-x_{i})\qquad (2)
\end{equation}


which in the continuum limit should approach an L' function normalized as.


\begin{equation}
\begin{array}{l l l l l l l}{{\displaystyle\int d x\,\rho(x)=1~}}&{{~~~~}}&{{~~~}}&{{~~~~~}}&{{~ }}}\end{array\end{array}
\end{equation}


In terms of this density function, one finds that.


\begin{equation}
{\cal F}\approx-\frac{9\pi}{8}N^{2+\alpha}\int d x d y\,\rho(x)\rho(y)\left(|x-y|+|x+y|\right)+\frac{\pi(8-N_{f})}{3}N^{1+\log\rho(x)\left|x\right|^{3}}\quad(4)
\end{equation}


where the large argument expansions have been used, and terms subleading in N have been
dropped. This only has non-trivial saddle points when both terms scale the same with N,
which demands that  = 1/2 and gives the famous result that F x N5/2. Extremizing the
free energy over normalized density functions then gives.


\begin{equation}
F\approx-{\frac{9\sqrt{2}\pi N^{5/2}}{5\sqrt{8-N_{f}}}}\qquad (5)
\end{equation}


This value of the free energy is to be identified with the renormalized on-shell action of the
supersymmetric AdSe solution. This identification yields the following relation between the
six-dimensional Newton's constant Gs and the parameters N and Ng of the dual SCFT,


\begin{equation}
G_{6}=\frac{5\pi\sqrt{8-N_{f}}}{27\sqrt{2}}\ N^{-5/2}\qquad (6)
\end{equation}

\end{document}
